\section{Discussion}
\label{sec:discussion}

Despite the rapid advancements in financial time-series generation, several fundamental challenges remain unresolved at the intersection of generative artificial intelligence and empirical financial econometrics. Bridging the gap between statistical expressiveness and market validity requires addressing core bottlenecks in causality, theoretical boundedness, cross-modal integration, and temporal non-stationarity.

\revaddblock{A first fundamental challenge concerns causal and invariant representation learning: current }FTSG models often learn spurious correlations specific to a particular market environment or training window. This stems from a lack of explicit disentanglement between true causal factors (e.g., shifts in fundamental central bank policies) and mere coincidental patterns (e.g., a specific high-frequency trading algorithm's temporary market dominance). Consequently, models that assume data are Independent and Identically Distributed (i.i.d.) \revdelblock{ fail drastically }\revaddblock{can degrade substantially }when encountering novel market regimes due to severe out-of-distribution (OOD) shifts \citep{arjovsky2019invariant, peters2016causal}. 
To address this, future FTSG frameworks must move toward environment-invariant learning. By treating specific market regimes (e.g., high vs. low volatility) as distinct training ``environments,'' techniques like Invariant Risk Minimization (IRM) \citep{arjovsky2019invariant} can be \revdelblock{ leveraged to eliminate }\revaddblock{used to reduce sensitivity to }spurious environmental correlations. Early attempts in forecasting demonstrate that enforcing representation invariance across regimes improves robustness \citep{cao2023irm}. Porting this paradigm to generative models \revdelblock{ will ensure that synthesized trajectories remain realistic even under unprecedented }\revaddblock{may improve the robustness of synthesized trajectories under novel }economic shifts.

\revdelblock{ 
A }\revaddblock{A second challenge is econometrically bounded generation. A }well-documented constraint of purely data-driven generative models (such as standard GANs\revdelblock{ and Diffusion Models}\revaddblock{~\citep{goodfellow2014gan} and diffusion models~\citep{ho2020ddpm,song2021score}}) is that they can replicate past empirical distributions without guaranteeing adherence to established market microstructure or stochastic process theory \citep{cont2001empirical}. For instance, unconstrained models might generate prices that breach arbitrage boundaries or fail to exhibit mean-reverting volatility.
Drawing inspiration from Physics-Informed Neural Networks (PINNs) in the physical sciences \citep{raissi2019physics}, future research should focus on \textit{econometrically-bounded} generation. Embedding economic priors---such as \revdelblock{ Geometric Brownian Motion }\revaddblock{geometric Brownian motion }(GBM)\revaddblock{~\citep{black1973pricing}}, stochastic volatility models, or jump-diffusion processes for extreme events~\citep{merton1976option}---directly into the generator's loss functions provides a powerful white-box constraint. Recent advancements have successfully integrated GBM priors into score-matching diffusion models \citep{tanaka2025cofindiff}. Expanding this to capture more complex stylized facts will ensure that synthesized sequences are not merely statistically plausible, but also economically rational.

\revaddblock{A third challenge is multi-modal time series generation. }Existing FTSG models broadly generate time-series data in isolation, neglecting the exogenous, multi-modal signals that drive actual market dynamics. However, financial markets are multi-causal systems: a macroeconomic announcement (textual metadata) triggers shifts in order flow (discrete events), which subsequently alter implied volatility surfaces (continuous manifolds). To capture this complexity, future work should pursue two promising avenues. First, cross-modal conditioned generation can encode broader heterogeneous data, such as central bank operations, limit order book states, and earnings call transcripts \cite{gholamrezaei2023conditioning, T2S, lozano2023dual}, using a unified representation space to precisely steer generative trajectories beyond narrow textual prompts \cite{Time-LLM, liu2025calf}. Second, joint multi-modal generation involves the simultaneous generation of linked data streams. By jointly modeling limit order books, financial news headlines, and derivative prices, generators can inherently learn cross-modal causal dependencies and lead-lag relationships \cite{lozano2023dual}, which is vital for simulating high-impact scenarios such as flash crashes.

\revaddblock{A fourth challenge is online adaptation under regime shifts. }Despite the zero-shot capabilities of modern time-series foundation models, they remain susceptible to concept drift---the frequent distributional shifts triggered by evolving economic cycles or regulatory changes. Therefore, static generative models will inevitably degrade over time. To avoid the notorious ``catastrophic forgetting'' associated with fine-tuning on new data \cite{mccloskey1989catastrophic, zenke2017continual}, FTSG must transition toward online adaptation.
One promising avenue is Bayesian Continuous Learning, utilizing Streaming Variational Bayes \cite{nguyen2017variational} to incrementally update the generator's latent space as new market observations arrive in real time. Alternatively, framing generation within an Online Reinforcement Learning framework \cite{moody2001learning}, where the model receives rewards based on the real-time fidelity of its synthetic samples relative to current market conditions, allows for near-instant calibration to abrupt regime shifts without retraining from scratch.

\revaddblock{Finally, we reflect on limitations, practical implications, and a concluding outlook. }While this survey synthesizes \revdelblock{ 130 }\revaddblock{131 }studies spanning 2021 to early 2026, several intrinsic limitations characterize the current empirical literature. Primarily, existing benchmarks remain disproportionately concentrated in liquid equity and cryptocurrency markets, with a comparative scarcity of standardized suites for fixed income, credit derivatives, and illiquid alternative assets. Furthermore, the absence of universally accepted, open-source matching engine simulators hampers standardized, apples-to-apples comparisons of downstream execution utility. 
Nonetheless, our proposed two-level taxonomy successfully unifies the field across task-oriented and technique-driven dimensions, providing a structured conceptual framework for future studies. By compiling curated databases, establishing a multi-dimensional evaluation protocol, and identifying theoretical boundaries, this review provides a \revdelblock{ definitive }\revaddblock{structured }foundation for researchers and practitioners developing next-generation financial generative models.